\begin{textsample}{POS Dim 2 – al – Score 22.00 – al/al\_em\_44.txt}  \label{ex:f2_pos_267}
Além da discussão e planejamento das unidades curriculares e formas do trabalho didático-pedagógico com as Áreas de Conhecimento e a BNCC, destaca -se a importância de planejar como acontece a rotina do : Aprofundamento, \textbf{Ofertas} \textbf{Eletivas} e Projeto de Vida, ou seja, a parte \textbf{flexível} do \textbf{currículo} do Ensino Médio.
Neste sentido, destacamos a seguir alguns aspectos fundamentais a serem considerados.
FORMAS DE ARTICULAÇÃO DAS UNIDADES CURRICULARES ( PROJETOS, ATELIÊS E LABORATÓRIOS DE ENSINO ) O planejamento das atividades escolares deve ocorrer de forma interdisciplinar, visando a integração e articulação das diferentes Áreas do Conhecimento, estudos e práticas e com foco nas habilidades a serem desenvolvidas.
A organização por Áreas do Conhecimento possibilita essa articulação dos conhecimentos de cada unidade curricular, mas preservando seus conceitos, procedimentos e atitudes.
Essa integração encontra espaço propício para acontecer no planejamento para o trabalho nos Ateliês Pedagógicos e Laboratórios de Aprendizagem, assim como no trabalho por meio de Projetos integradores e desenvolvimento de \textbf{oficinas} mão na massa.
As unidades curriculares devem estar sempre articuladas de forma a potencializar as aprendizagens.
Na elaboração dos \textbf{Itinerários} \textbf{Formativos} de Áreas do Conhecimento é essencial assegurar o desenvolvimento dos quatro eixos \textbf{estruturantes} e suas respectivas habilidades, bem como certificar -se de que os objetos de conhecimento abordados não sejam os mesmos da Formação Geral, mas, de fato, aprofundem ou ampliem aqueles já trabalhados.
Da mesma forma, os \textbf{Itinerários} \textbf{Formativos} de Formação \textbf{Técnica} e Profissional precisam levar em consideração os quatro eixos \textbf{estruturantes} e suas habilidades no Módulo de Formação para o Mundo do Trabalho, além de garantir que as Matrizes Curriculares dos \textbf{Cursos} \textbf{Técnicos} dialoguem com as Competências Gerais da BNCC.
Para os \textbf{Itinerários} \textbf{Formativos}, orienta -se que, no planejamento, a escola ou \textbf{rede} de ensino possa responder às seguintes situações : Quais as possibilidades de Aprofundamento?
Os Aprofundamentos representam um diferencial no \textbf{currículo} do Ensino Médio.
Para que as aprendizagens promovidas pela Formação Geral sejam aprofundadas e ampliadas, é necessário que os \textbf{Itinerários} \textbf{Formativos} de Áreas do Conhecimento abordem temáticas contemporâneas contextualizadas e que levem em consideração os interesses dos estudantes.
No caso do Itinerário \textbf{Formativo} de Formação \textbf{Técnica} e Profissional, o aprofundamento ocorre a partir do desenvolvimento de habilidades básicas exigidas pelo mundo do trabalho, e habilidades específicas dos \textbf{Cursos} \textbf{Técnicos}, \textbf{Cursos} de \textbf{Qualificação} \textbf{Profissional} ( FICs ) ou Programa de Aprendizagem Profissional escolhidos.
Quais as possibilidades do Projeto de Vida?
O Projeto de Vida deve ser uma unidade curricular, presente nos três anos do Ensino Médio, com \textbf{carga} \textbf{horária} mínima de dois tempos de aula por semana e material específico para trabalhar didático-pedagógico e processos de ensino e aprendizagem.
O trabalho com o projeto de vida deve estar \textbf{alinhado} com a comunidade escolar e ocorrer também em diversos momentos na escola, não apenas no tempo de aula do componente curricular.
Isso possibilitará o desenvolvimento do projeto de vida de cada estudante, estimulando o autoconhecimento, a ampliação da sua percepção de mundo, preparando o jovem para fazer planejar e fazer escolhas em toda a sua \textbf{trajetória}.
O Projeto de Vida tem caráter eminentemente transversal à \textbf{trajetória} do estudante do ensino médio, bem como à \textbf{trajetória} já construída na educação básica.
Neste sentido, todas as propostas de ensino presentes no \textbf{itinerário} \textbf{formativo} e formação geral, devem considerar grande articulação e integração com a unidade curricular Projeto de Vida.
Tudo \textbf{concorre} ou converge para o apoio e desenho do Projeto de Vida do estudante.
Todo o processo deve ser permeado por vivências que lhes permitam desenvolver competências como autoconfiança, determinação e resiliência, dentre outros.
Quais as possibilidades da \textbf{Oferta} \textbf{Eletiva}?
As \textbf{eletivas} devem ser idealizadas com temas que promovam diferentes vivências e aprendizagens, enriquecendo o Itinerário \textbf{Formativo}.
Devem estar articuladas com as Áreas do Conhecimento, \textbf{itinerários} \textbf{formativos}, os eixos \textbf{estruturantes} e as Competências Gerais da BNCC, ter caráter lúdico e prático e, sempre com intencionalidade pedagógica.
A recomendação é que as \textbf{eletivas} sejam construídas pelos professores, considerando territórios, escutas qualificadas, sugestões e interesses dos estudantes, mediante um processo validado pela escola e, no caso da \textbf{rede} \textbf{estadual} de ensino, pela SEDUC/AL.
Orienta -se a composição de uma \textbf{catálogo} de \textbf{eletivas}, o qual será formado pelas unidades curriculares \textbf{eletivas} validadas, assim podem passar a compor um \textbf{catálogo} de possibilidades de \textbf{oferta} para a \textbf{rede} ou escola.
Destaca -se a importância de escutas dos estudantes e mundo do trabalho, promovendo observação qualificada do ambiente, para definir quais os melhores caminhos para preparação e intervenção sustentável.
Na Formação \textbf{Técnica} e Profissional, as FICs ( \textbf{Curso} de \textbf{Qualificação} \textbf{Profissional} ) também podem ser \textbf{ofertadas} como \textbf{eletivas}.
A \textbf{carga} \textbf{horária} do Itinerário \textbf{Formativo}, ao longo dos três anos de formação estudante, deve ser distribuída, conforme Aprofundamento, Projeto de Vida e \textbf{Eletivas}, que devem ser propostas pedagógicas da \textbf{rede} de ensino e unidade de ensino \textbf{ofertante}.
Por fim, é importante observar que o desenvolvimento das habilidades previstas nos \textbf{Itinerários} \textbf{Formativos} demandam a aplicação de metodologias ativas e diversificadas, bem como o foco deve ser maior no protagonismo dos estudantes.
Tais especificidades também requerem criatividade e inovação no que diz respeito à elaboração de arranjos curriculares que oportunizem mudanças significativas na organização de tempos, espaços e práticas escolares, de forma a assegurar que façam mais sentido e gerem mais aprendizagem e desenvolvimento para jovens do século XXI.

% matched lemmas: alinhar, carga, catálogo, concorrer, currículo, curso, eletivo, estadual, estruturante, flexível, formativo, horário, itinerário, oferta, ofertante, ofertar, oficina, profissional, qualificação, rede, trajetória, técnico
\end{textsample}
